\documentclass[10pt,a4paper]{amsart}
\usepackage[margin=19mm]{geometry}
\usepackage{amsmath,amssymb,mathtools}
\usepackage[scr=rsfs,cal=euler]{mathalfa}
\usepackage{xfrac}
\usepackage[unicode,hidelinks,bookmarks=false]{hyperref}
\newcommand{\ie}{i.e.\ }
\newcommand{\cf}{cf.\ }
\newcommand{\resp}{respectively\ }
\newcommand{\Subsection}{\subsection}
\providecommand{\nobreakdash}{\nobreak\hbox{-}}
\setlength{\emergencystretch}{2em}
\multlinegap0pt
\usepackage{amssymb}
\usepackage{mathtools}
\usepackage[shortlabels]{enumitem}
\newtheorem{lemma}{Lemma}
\newtheorem{proposition}{Proposition}
\newtheorem{theorem}{Theorem}
\newcommand{\Q}{\mathbb{Q}}
\newcommand{\Tr}{\mathrm{Tr}}
\newcommand{\Z}{\mathbb{Z}}
\newcommand{\cO}{\mathcal{O}}
\newcommand{\diag}{\mathrm{diag}}
\newcommand{\sh}{\mathrm{sh}}
\title{Minkowski shapes of pure number fields}
\author{Khai-Hoan Nguyen-Dang}
\date{Source: arXiv:2606.07956v1 (CC BY 4.0)}
\begin{document}
\maketitle

\noindent\emph{Source and licence.} Minkowski shapes of pure number fields. Version arXiv:2606.07956v1 (CC BY 4.0), distributed by arXiv under the Creative Commons Attribution 4.0 International licence, http://creativecommons.org/licenses/by/4.0/, as recorded in the arXiv licence metadata for this exact version and shown on the abstract page https://arxiv.org/abs/2606.07956v1. Full licence notice: \texttt{licenses/arxiv-2606.07956-CC-BY-4.0.txt}.\par\smallskip

\noindent\textit{Editorial scope.} Counted unit: the article's theorem thm:even-Sh-ambiguity (source0.tex lines 2198--2206). The article's complete original proof of it is delivered with it (source0.tex lines 2208--2357, one \texttt{proof} environment of 150 lines). No mathematics is written, repaired or re-derived: the statement and the proof are the article's own lines, in the article's own notation, and no retained line is substituted or re-flowed.  Retained uncounted as the source-local context the proof needs: lines 560--572; lines 978--986; lines 2167--2182.  Nothing was omitted from any retained span, so the package carries no omission record.  No retained line contains a citation, so the wrapper carries no bibliography and no bibliography entry.  No retained line contains a figure environment, an image-inclusion command or drawing code, so no image is delivered and the package declares no asset dependency.  Editorial formatting only, in addition: an AMS \texttt{amsart} preamble that restates the article's own declarations and macros used by the retained lines, namely the environments \texttt{lemma}, \texttt{proposition}, \texttt{theorem} on separate counters, together with the standard packages those lines need.  The retained environments print in this document as Lemma 1, Proposition 1, Theorem 1 in order of appearance, whereas the article prints the same items with the numbering of its own full document.  The complete native source of the article as distributed in the arXiv e-print for this exact version is bundled unchanged as \texttt{sources/202292/source0.tex}.
\begin{lemma}[Exact sequence and rank]\label{lem:exact-sequence}
Let $K/\Q$ have degree $n$, and let $\mathcal L\subset \cO_K$ be a full-rank $\Z$-submodule containing $1$.
Then the map
\[
\mathcal L\longrightarrow \mathcal L^\perp,\qquad \alpha\longmapsto \alpha^\perp
\]
is a surjective homomorphism with kernel $\Z\cdot 1$.
Consequently, $\mathcal L^\perp$ is a free $\Z$-module of rank $n-1$ and
\[
0\longrightarrow \Z\cdot 1\longrightarrow \mathcal L\longrightarrow \mathcal L^\perp\longrightarrow 0
\]
is exact.
\end{lemma}

\begin{proposition}[Orthogonality of monomials]\label{prop:orthogonality}
Let $a$ be admissible and let $1\le i,j\le n-1$.
Then
\[
\langle \theta^i,\theta^j\rangle_\infty=0\quad\text{if }i\neq j,
\qquad
\langle \theta^m,\theta^m\rangle_\infty=n\,|a|^{2m/n}\quad(1\le m\le n-1).
\]
\end{proposition}

\begin{theorem}[Nguyen-Dang--Nguyen \(\alpha\)-monogeneity criterion]
\label{thm:NDH-monogenic}
Let \(K=\Q(\alpha)\) with \(\alpha\) a root of an irreducible polynomial
\[
x^n-m\in \Z[x].
\]
Then
\[
\cO_K=\Z[\alpha]
\]
if and only if \(m\) is squarefree and
\[
v_p(m^p-m)=1
\qquad\text{for every prime }p\mid n.
\]
\end{theorem}

\begin{theorem}[Sign ambiguity for the shape]\label{thm:even-Sh-ambiguity}
Assume that $n=2r \geq 4$ is even.
Then there exist infinitely many squarefree integers $a>0$ such that
\[
\operatorname{sh}(K_a)=\operatorname{sh}(K_{-a})
\qquad\text{but}\qquad
K_a\not\simeq K_{-a}.
\]
\end{theorem}

\begin{proof}
Set
\[
N:=4\prod_{\substack{p\mid n\\ p\ \mathrm{odd}}} p^2.
\]
For each odd prime \(p\mid n\), choose a residue class \(c_p\pmod{p^2}\) satisfying
\[
2c_p\equiv 1+p \pmod{p^2}.
\]
Also impose the congruence
\[
c\equiv 1 \pmod 4.
\]
By the Chinese remainder theorem, there exists an integer \(c\) modulo \(N\) satisfying
\[
c\equiv c_p \pmod{p^2}\quad (p\mid n,\ p\ \mathrm{odd}),
\qquad
c\equiv 1\pmod 4.
\]
Since \(c\) is coprime to \(N\), Dirichlet's theorem yields infinitely many primes
\[
q\equiv c \pmod N.
\]
Fix such a prime \(q\) and put
\[
a:=2q.
\]

Then \(a>0\) and \(a\) is squarefree. Hence \(a\) is \(n\)th-power-free, and for every prime
\(p\mid n\) we have \(v_p(a)\in\{0,1\}\); therefore Hypothesis~\textup{(H)} holds.
Moreover, since \(q\mid a\) but \(q^2\nmid a\), both
\[
x^n-a\qquad\text{and}\qquad x^n+a
\]
are Eisenstein at \(q\), hence irreducible over \(\Q\). Thus \(a\) and \(-a\) are admissible.

We now verify the criterion.

If \(p=2\), then \(a=2q\) with \(q\) odd, so
\[
v_2(a^2-a)=v_2\bigl(2q(2q-1)\bigr)=1,
\]
and
\[
v_2((-a)^2-(-a))=v_2(a^2+a)=v_2\bigl(2q(2q+1)\bigr)=1.
\]

Now let \(p\mid n\) be odd. Since \(q\equiv c\pmod{p^2}\), we have
\[
a=2q\equiv 2c\equiv 1+p \pmod{p^2}.
\]
In particular \(p\nmid a\), and
\[
a^{p-1}\equiv (1+p)^{p-1}\equiv 1+(p-1)p \equiv 1-p \not\equiv 1 \pmod{p^2}.
\]
Since \(p\nmid a\), Fermat's little theorem gives
\[
a^{p-1}\equiv 1 \pmod p,
\]
while the displayed congruence shows
\[
a^{p-1}\not\equiv 1 \pmod{p^2}.
\]
Hence
\[
v_p(a^{p-1}-1)=1,
\]
and therefore
\[
v_p(a^p-a)=v_p(a)+v_p(a^{p-1}-1)=1.
\]
Since \(p\) is odd,
\[
(-a)^p-(-a)=-(a^p-a),
\]
hence also
\[
v_p((-a)^p-(-a))=1.
\]

Thus Theorem~\ref{thm:NDH-monogenic} gives
\[
\cO_{K_a}=\Z[\theta],\qquad \theta^n=a,
\]
and
\[
\cO_{K_{-a}}=\Z[\vartheta],\qquad \vartheta^n=-a.
\]

Because \(a\) is squarefree, the strong decomposition of \(a\) has
\[
a_1=a,\qquad a_j=1\quad (2\le j\le n-1),
\]
and the strong decomposition of \(-a\) has the same positive part. Hence
\[
C_m(a)=C_m(-a)=1
\qquad (1\le m\le n-1).
\]
Therefore the normalized monomials in the two fields are simply
\[
e_m(a)=\theta^m,
\qquad
e_m(-a)=\vartheta^m
\qquad (1\le m\le n-1).
\]

Since \(\Tr(\theta^m)=\Tr(\vartheta^m)=0\) for \(1\le m\le n-1\), Lemma~\ref{lem:exact-sequence}
shows that
\[
(n\theta,\dots,n\theta^{n-1})
\]
is a \(\Z\)-basis of \(\cO_{K_a}^{\perp}\), and
\[
(n\vartheta,\dots,n\vartheta^{n-1})
\]
is a \(\Z\)-basis of \(\cO_{K_{-a}}^{\perp}\).

By Proposition~\ref{prop:orthogonality}, for \(1\le i,j\le n-1\),
\[
\bigl\langle J(\theta^i),J(\theta^j)\bigr\rangle_\infty
=
\begin{cases}
0,& i\neq j,\\[2mm]
n\,a^{2i/n},& i=j,
\end{cases}
\]
and exactly the same formula holds with \(\vartheta\) in place of \(\theta\), because only \(|a|\)
enters. Hence the Gram matrices of the above trace-zero bases are identical:
\[
\mathrm{Gram}\bigl(J(n\theta),\dots,J(n\theta^{n-1})\bigr)
=
\mathrm{Gram}\bigl(J(n\vartheta),\dots,J(n\vartheta^{n-1})\bigr)
=
n^3\diag\!\bigl(a^{2/n},a^{4/n},\dots,a^{2(n-1)/n}\bigr).
\]
Therefore
\[
\sh(K_a)=\sh(K_{-a}).
\]

Finally, since \(n=2r\) is even and \(a>0\), the polynomial \(x^n-a\) has exactly two real roots,
whereas \(x^n+a\) has no real roots. Thus
\[
\operatorname{sig}(K_a)=(2,r-1),
\qquad
\operatorname{sig}(K_{-a})=(0,r),
\]
so \(K_a\not\simeq K_{-a}\).
Since there are infinitely many such primes \(q\), the theorem follows.
\end{proof}

\end{document}
